\documentclass[11pt,a4paper]{article}
\usepackage[margin=25mm]{geometry}
\usepackage{amsmath,amssymb,amsthm,hyperref}
\setlength{\emergencystretch}{2em}
\newtheorem{lemma}{Lemma}\newtheorem{theorem}[lemma]{Theorem}
\newcommand{\gen}[1]{\langle#1\rangle}
\newcommand{\Sym}{\operatorname{Sym}}
\newcommand{\Alt}{\operatorname{Alt}}
\title{The ten-point symmetric factor}
\author{Human Compression project: a component of the L019 recovery}
\date{}
\begin{document}\maketitle
All actions below are actual permutation actions. A marking is
\emph{fractured} when every maximal omission is intransitive.
An endpoint anchor is moved by the first mark and fixed by all
later marks, or by the reversed version of that condition.

\begin{lemma}\label{orbits}
In a fractured transitive marking, each mark of every prefix is
essential to that prefix's orbit partition.
\end{lemma}
\begin{proof}
The orbits of a generated permutation group are the equivalence
classes generated by the pairs $x,x^s$ for its generators $s$.
If a prefix and one of its omissions have the same orbit partition,
adjoining the same next generator to both gives the same partition
again. Inductively they still have the same partition after all
remaining generators have been adjoined. The full marking is
transitive, so the corresponding full omission would be transitive
too, contradicting fracture.
\end{proof}

\begin{theorem}\label{ten}
Let $(t_0,\ldots,t_5)$ be a fractured natural $S_{10}$ string
C-marking with endpoint anchor $a$. Suppose that, on the same
ten-point domain, it satisfies
\[
 P_j\cap\Alt(X_{j-1})\le P_{j-1}\quad(1\le j\le6),
 \qquad P_j=\gen{t_0,\ldots,t_{j-1}},\quad X_j=a^{P_j}.
                                                        \tag{1}
\]
Then $t_0$ is a pure transposition.
\end{theorem}
\begin{proof}
We give a complete finite exclusion of a nonpure $t_0$.
Conjugate the anchor to point $1$. All involutions of $S_{10}$
are obtained by taking the five conjugacy classes represented by
products of one, two, three, four or five disjoint transpositions.
There are $9495$ nonidentity involutions, of which $2619$ fix $1$.
At the first position we require movement of $1$ and support
larger than two. At each later position we try every involution
fixing $1$ which commutes with every previous mark except its
immediate predecessor, does not commute with that predecessor,
and does not belong to the group of the previous prefix.

Adjacent noncommutation is necessary here. In a transitive anchored
string action, let $H$ be generated by the connected component of
the noncommutation diagram containing the first mark. Every other
component centralizes $H$ and fixes $1$. The full orbit of $1$
therefore equals its $H$-orbit, so $H$ is transitive. A permutation
centralizing a transitive action and fixing one point fixes all
points. There can consequently be no nonidentity mark in another
component. Membership in the previous prefix is excluded by IP.

Only conjugacy repetitions are removed. Initially conjugate by
$(S_{10})_1$. After marks have been chosen, conjugate the next
candidate only by the subgroup fixing $1$ and centralizing every
chosen mark. This subgroup preserves all stated conditions.
Each of its orbits contributes one representative, so induction
on the prefix length proves that every possible marking remains
represented. No arbitrary graph identification or restriction
quotient is involved.

For a prefix group $H$ of length $l$, the necessary order condition
is $[S_{10}:H]\ge2^{6-l}$: each remaining independent generator
at least doubles its preceding subgroup. Every trial prefix is
also tested by Lemma~\ref{orbits} and by (1).
The latter is computed as an actual intersection with the group
of even permutations supported on the previous anchor orbit.
For orbits of size at most two that group is trivial.

The interval intersection tests are also necessary consequences
of the original IP. When appending a mark at position $l$, put
\[
 H_i=\gen{t_i,\ldots,t_{l-1}},\qquad
 J_i=\gen{t_i,\ldots,t_l}.
\]
For every $0\le i\le l-2$ require
\[
                         H_i\cap J_{i+1}=H_{i+1}.       \tag{2}
\]
The old groups $H_i$ are retained from the preceding step.
These tests avoid repeatedly constructing every subset group at
every node. For the exclusion, only their necessity is needed.
At every terminal node the implementation nevertheless checks
all pairs of all $64$ generated subset groups directly. Thus its
reported terminal markings also satisfy the full IP.

The complete counts of visited prefixes, including the empty one,
are
\[
                      1,4,197,1775,645,208,18.
\]
All four nonpure first-mark types are exhausted. The orbit test
rejects $3949$ candidates, (2) rejects $1336$, and (1) rejects
$1402$ further candidates, in that order. None of the eighteen
terminal groups has order $10!=3628800$. Since all are subgroups
of the given faithful $S_{10}$, none is the required natural group.
The literal source program, full terminal receipt, terminal orders
and orbit sizes are retained. Every rejection was justified by a
necessary condition, and every conjugacy reduction was exhaustive.
This proves the finite exclusion and hence the theorem.
\end{proof}

\begin{theorem}\label{application}
Consider the actual critical one-point configuration retained in
the L019 recovery: a fractured natural symmetric factor $P$ of
even degree $m$ and rank $m/2+1$, and an even six-point rank-three
$A_5$ factor $E$, supported on $X,Y$ with $X\cap Y=\{a\}$.
Their labels are consecutive disjoint intervals, and $a$ is an
endpoint anchor in each. Assume the whole marking is a string
C-marking. If $m=10$, the whole marking has an
internal pure perfect split. Together with the acquired smaller
cases, only $m\ge12$ remains in this critical configuration.
\end{theorem}
\begin{proof}
Orient $P$ with its anchor first. The acquired odd-attachment
theorem supplies (1) from the IP of the whole marking and the
actual external transitive group $E$. It requires no odd element
in $E$ for this even-kernel conclusion. Theorem~\ref{ten} gives
$t_0=(a\ c)$.

The tail of $P$ fixes $a$ and is transitive on $X\setminus\{a\}$.
Indeed, if it had more than one orbit there, adjoining the sole
edge $ac$ could join $a$ to only one of them, contrary to
transitivity of $P$. In the whole marking, omission of $t_0$
therefore leaves precisely $Y$ and $X\setminus\{a\}$ as its
orbits. The two label intervals are supported and transitive on
their respective sets, and the omitted mark is pure. This is an
internal pure perfect split. The acquired pure-split criterion
supplies its valid rank and degree extension. The separately
reviewed small-symmetric component settles $m\le8$; no larger
degree is being inferred from the present finite calculation.
\end{proof}

\paragraph{Attribution and limits.}
The principal source is Cameron--Fernandes--Leemans,
\emph{The number of string C-groups of high rank},
arXiv:2212.12723v2. The critical reduction, supported even-kernel
theorem and precise split criterion are acquired with their
attributions and review chains. The present contribution is the
orbit-partition pruning argument and this complete bounded
ten-point exclusion. Permutation-group operations use GAP; the
conjugacy and branching completeness argument is given above.
An independent small-group check compares the interval tests with
all subset intersections, and a retained pure-first $S_{10}$
example is checked as a positive control.

Two earlier degree-ten searches stopped at their declared caps
without a certificate. Their receipts remain unchanged. A first
attempt at the present program stopped before search because its
global variable name was reserved in GAP; the corrected complete
run and the precise change are recorded. These earlier executions
carry no exclusion credit. The general critical case, contraction,
the stability correspondence under isomorphism and duality, and
the exact base counts remain pending. No whole-original, new
question, Lean or compression-measurement credit is asserted.
\end{document}
